\documentclass[12pt,a4paper]{report}
\usepackage[utf8]{inputenc}
\usepackage{amsmath,amssymb,amsfonts}
\usepackage{geometry}
\geometry{top=1in, bottom=1in, left=1.25in, right=1in}
\usepackage{graphicx}
\usepackage{booktabs}
\usepackage{array}
\usepackage{hyperref}
\usepackage{setspace}
\onehalfspacing

\begin{document}
\begin{titlepage}
    \centering
    \vspace*{0.5cm}
    {\Large \textbf{NATIONAL INSTITUTE OF TECHNOLOGY HAMIRPUR}}\\[0.3cm]
    {\large DEPARTMENT OF MECHANICAL ENGINEERING}\\[2.0cm]
    {\LARGE \textbf{OPTIMIZATION OF MACHINING PARAMETERS IN CNC MILLING USING RESPONSE SURFACE METHODOLOGY (RSM)}}\\[2.5cm]
    \textbf{A Project Report Submitted in Partial Fulfillment of the Requirements for the Degree of}\\[0.3cm]
    \textbf{Bachelor of Technology in Mechanical Engineering}\\[2.0cm]
    \textbf{Submitted by:}\\[0.2cm]
    \textbf{Banoth Chandu Nayak}\\[3.0cm]
    \textbf{MAY 2026}
\end{titlepage}

\chapter*{Abstract}
This project investigates the optimization of machining parameters in dry peripheral end-milling of structural alloy steel (AISI 4340, 32 HRC) using TiAlN-coated solid carbide end-mills. A 3-factor, 20-run Central Composite Design (CCD) was implemented to evaluate spindle speed (2000--3000 RPM), feed rate (200--500 mm/min), and axial depth of cut (0.75--1.75 mm). Second-order quadratic response surface models were fitted for surface roughness ($R_a$) and tool flank wear ($V_b$). Constrained numerical optimization yielded optimal cutting settings of $N = 2800\text{ RPM}$, $f = 250\text{ mm/min}$, and $a_p = 1.05\text{ mm}$, successfully validating a \textbf{14.01\% reduction in surface roughness} relative to the central baseline.

\end{document}
